\subsection{System Check and Installation}
\label{subsec:systempruefung-installation}

The read-only \texttt{doctor} checks runtimes, configuration, paths,
dependencies, and ports according to the profile. Disabled layers are not an
error. In general, missing Docker produces only a warning; container,
database, and Tauri prerequisites have their own stricter diagnostic paths
\parencite{kleiveist2026tooling,kleiveist2026configuration}.

\texttt{install} sets up only active components: frontend dependencies are
preferentially installed from the lockfile, and backend dependencies are
installed in a virtual environment. The dedicated \path{tools/.venv} is set up
for repository tools, including Ruff, regardless of profile; the
\path{backend/.venv} remains separate. Database packages and Chromium are
added only when required; for Python, a pip/venv fallback is available
alongside \texttt{uv}. A local \texttt{.env} remains unchanged. \texttt{tauri
install} handles native operating-system and Rust prerequisites separately.
Specifically, the frontend uses \texttt{npm ci} when a lockfile is present. An
active backend receives its own environment and only the profile-dependent
requirement sets. Chromium is installed only when Playwright is configured;
disabled layers are skipped.

At the historical observation state \texttt{461fc751}, generated backend
profiles had an installation discrepancy relevant to governance: \texttt{install}
could provide Ruff through \path{backend/.venv} and omit
\path{tools/.venv}, while the quality adapter expected Ruff there or on the
global search path. In a freshly generated \texttt{web-cloud} profile,
installation completed successfully, but the subsequent quality command
reported Ruff as unavailable. Core CI avoided this combination through
\texttt{install --skip-backend}; the generated backend-profile jobs did not.
This commit-bound observation explains the profile failures at the time and
remains part of the historical evidence
\parencite{kleiveist2026quality-implementation,kleiveist2026quality-ci}.

In the v1.0.0 release state, the discrepancy is corrected: \texttt{install}
provides \path{tools/.venv} regardless of profile, and the quality adapter uses
this dedicated tooling runtime. The local \texttt{install} $\rightarrow$
\texttt{quality} path therefore has a uniform contract for Ruff resolution.
This later correction does not rewrite the commit-bound CI results of the
baseline.

The README requires Git, Python 3.11 or later, Node.js 24 or later, and Rust
Stable for desktop profiles. The general Doctor does not check Git or these
version ranges; the Tauri Doctor accepts any active Rust toolchain. CI, by
contrast, explicitly sets up the specified major versions of Python and Node.
Local enforcement therefore differs from the documentation
\parencite{kleiveist2026governance-snapshot,kleiveist2026ci}.
This limitation concerns how the prerequisite is checked, not whether the
documented prerequisite applies.
\beginnerinput{workspace/statement/500_tooling_workflow/530}
